%%%%%%%% SOSP 2026 PAPER TEMPLATE %%%%%%%%%%%%%%%%%
%
% The 32nd ACM Symposium on Operating Systems Principles
% September 30, 2026
%
% Format: ACM SIGPLAN, <= 12 pages technical content (excluding references)
%         A4 or US letter, 178x229mm (7x9") text block
%         Two-column, 8mm separation, 10pt on 12pt leading
%
% Official CFP: https://sigops.org/s/conferences/sosp/2026/cfp.html
% ACM Template: https://www.acm.org/publications/proceedings-template
%
% IMPORTANT NOTES:
% - Double-blind review (use paper ID, not author names)
% - Anonymized system/project name required (different from arXiv/talks)
% - Optional Artifact Evaluation after acceptance
% - Author response period available:
%   --> LIMITED TO: correcting factual errors + addressing questions
%   --> NO new experiments or additional work
%   --> Keep under 500 words
% - Supplementary material allowed (reviewers not required to read)
% - Figures/tables readable without magnification, color encouraged
% - Pages numbered, references hyperlinked
%
% WHAT SOSP VALUES:
% - Groundbreaking work in significant new directions
% - Significant problem motivation
% - Interesting, compelling solution with demonstrated practicality
% - Clear contributions and advances beyond previous work
% - Papers addressing new problems may be evaluated differently
%   from those in established areas

\documentclass[sigplan,10pt]{acmart}

% Remove copyright/permission footer for submission
\renewcommand\footnotetextcopyrightpermission[1]{}
\settopmatter{printfolios=true}

% Remove ACM reference format for submission
\setcopyright{none}
\renewcommand\acmConference[4]{}
\acmDOI{}
\acmISBN{}

% Recommended packages for systems papers
\usepackage{booktabs}       % Professional tables
\usepackage{xspace}
\usepackage{subcaption}     % Side-by-side figures
\usepackage{algorithm}      % Algorithm environment
\usepackage{algorithmic}    % Pseudocode formatting
\usepackage{listings}       % Code listings
\usepackage[capitalize,noabbrev]{cleveref}  % Smart cross-references

% Code listing style
\lstset{
  basicstyle=\footnotesize\ttfamily,
  numbers=left,
  numberstyle=\tiny,
  xleftmargin=2em,
  breaklines=true,
  tabsize=2,
  showstringspaces=false,
  frame=single,
  captionpos=b,
  language=C  % Default language; change as needed
}

% Custom commands -- replace \system with your anonymized name
\newcommand{\system}{SystemName\xspace}
\newcommand{\eg}{e.g.,\xspace}
\newcommand{\ie}{i.e.,\xspace}
\newcommand{\etal}{\textit{et al.}\xspace}
\newcommand{\para}[1]{\smallskip\noindent\textbf{#1.}}
\newcommand{\parait}[1]{\smallskip\noindent\textit{#1.}}

% Systems-specific unit macros
\newcommand{\us}{\,$\mu$s\xspace}
\newcommand{\ms}{\,ms\xspace}
\newcommand{\ns}{\,ns\xspace}
\newcommand{\GB}{\,GB\xspace}
\newcommand{\MB}{\,MB\xspace}
\newcommand{\TB}{\,TB\xspace}
\newcommand{\Gbps}{\,Gbps\xspace}

\begin{document}

\title{Your Paper Title Here}

% Anonymized for submission -- use paper ID
\author{Paper \#XXX}
\affiliation{%
  \institution{Anonymous}
  \country{}}

% Camera-ready (uncomment and fill in):
% \author{Author One}
% \affiliation{%
%   \institution{University/Company}
%   \city{City}
%   \country{Country}}
% \email{email@example.com}
%
% \author{Author Two}
% \affiliation{%
%   \institution{University/Company}
%   \city{City}
%   \country{Country}}
% \email{email@example.com}

\begin{abstract}
% Guidelines for a strong SOSP abstract:
% - State the OS/systems principle advanced
% - Identify why existing approaches are insufficient
% - Describe your approach and key insight
% - Quantify with concrete numbers
%
% Keep to 150--200 words. SOSP values groundbreaking contributions
% to operating systems principles.

We present \system, a [describe system] that [key capability] for
[OS/systems problem].
[Problem: why existing OS approaches fall short.]
Our key insight is that [fundamental observation about systems design].
\system realizes this insight through [N] novel mechanisms:
(1)~[technique A] and (2)~[technique B].
We evaluate \system on [workloads] and demonstrate [X]$\times$
improvement in [metric] over [baseline], while maintaining
[reliability/consistency/other property].
\end{abstract}

\maketitle
\pagestyle{plain}

%----------------------------------------------------------------------
\section{Introduction}
\label{sec:intro}

% SOSP values groundbreaking work. Structure your introduction to show:
% 1. Important problem in systems principles (1--2 paragraphs)
% 2. Fundamental limitation of existing approaches (1 paragraph)
% 3. Key insight -- a new principle or observation (1 paragraph)
% 4. System design and approach overview (1 paragraph)
% 5. Contributions (bulleted list)
% 6. Results preview with concrete numbers (1 paragraph)
%
% SOSP encourages papers that open significant new directions.
% Evaluation criteria for papers addressing new problems may differ
% from those in established areas.

Operating systems must [challenge] as modern hardware and workloads
evolve~\cite{ghemawat2003gfs}. The traditional approach of [prior method]
was designed for [assumptions], but [new trend] fundamentally changes
the landscape~\cite{corbett2013spanner}.

\para{Fundamental Limitation}
Existing systems~\cite{decandia2007dynamo, hunt2010zookeeper}
rely on [assumption]. We show that this assumption breaks down
when [condition], leading to [consequence].

\para{Key Insight}
We observe that [fundamental systems principle/observation]. This
insight enables a new approach where [high-level description].

\para{\system Overview}
Building on this insight, we design \system, which introduces:
(1)~[mechanism A] for [purpose], and (2)~[mechanism B] for [purpose].
Together, these enable [combined capability].

We make the following contributions:
\begin{itemize}
  \item We identify a fundamental limitation in [existing approach]
        and formalize the problem (\cref{sec:background}).
  \item We design \system with [N] novel mechanisms: [list]
        (\cref{sec:design}).
  \item We prove that \system provides [formal guarantee] under
        [conditions] (\cref{sec:correctness}).
  \item We implement \system in [X]K lines of [language] and evaluate
        on [workloads], demonstrating [X]$\times$ improvement
        (\cref{sec:evaluation}).
\end{itemize}

%----------------------------------------------------------------------
\section{Background and Motivation}
\label{sec:background}

\subsection{System Model and Assumptions}

Define your system model, including the hardware, software, and
failure assumptions.

\subsection{Limitations of Existing Approaches}

Explain why current systems are insufficient. Use concrete examples.

\subsection{Motivating Measurements}

% Real measurements on production systems or realistic workloads
% are highly valued at SOSP.
\Cref{fig:motivation} shows [measurement] that illustrates the
fundamental problem.

\begin{figure}[t]
  \centering
  \fbox{\parbox{0.9\columnwidth}{\centering\vspace{3em}
    \textit{Measurement study showing the fundamental problem: \\
    e.g., latency distribution, throughput breakdown, or failure analysis}
  \vspace{3em}}}
  \caption{[Description.] We analyze [N] hours of production workload
    and find that [X]\% of [operations] violate [property],
    motivating [your approach].}
  \label{fig:motivation}
\end{figure}

%----------------------------------------------------------------------
\section{Design}
\label{sec:design}

% Present your system design clearly and rigorously.
% SOSP papers often include formal guarantees or invariants.

\Cref{fig:architecture} presents the architecture of \system.

\begin{figure*}[t]
  \centering
  \fbox{\parbox{0.9\textwidth}{\centering\vspace{4em}
    \textit{System architecture diagram showing components, \\
    data flow, and control flow}
  \vspace{4em}}}
  \caption{Architecture of \system. [Describe the components
    and their interactions.]}
  \label{fig:architecture}
\end{figure*}

\subsection{[Mechanism A]}
\label{sec:design:a}

Describe the first key mechanism. Its core logic is specified
in \cref{alg:mechanism}.

\begin{algorithm}[t]
  \caption{[Mechanism A] in \system}
  \label{alg:mechanism}
  \begin{algorithmic}[1]
    \STATE \textbf{Input:} request $r$, state $S$, configuration $C$
    \STATE \textbf{Output:} response and updated state
    \STATE \textbf{Invariant:} $\forall t: \text{Consistent}(S_t)$
    \STATE Acquire lock on $S.\text{partition}(r.\text{key})$
    \IF{$r.\text{type} = \text{READ}$}
      \STATE $v \leftarrow S.\text{Get}(r.\text{key}, r.\text{timestamp})$
      \STATE \textbf{return} $(v, S)$
    \ELSE
      \STATE $S' \leftarrow S.\text{Apply}(r.\text{mutation})$
      \STATE Replicate $S'$ to $C.\text{replicas}$
      \STATE Wait for quorum acknowledgment
      \STATE \textbf{return} $(\text{OK}, S')$
    \ENDIF
  \end{algorithmic}
\end{algorithm}

\subsection{[Mechanism B]}
\label{sec:design:b}

Describe the second key mechanism. The consistency guarantee can
be formally expressed as:
\begin{equation}
  \label{eq:consistency}
  \forall r_1, r_2 \in \mathcal{R}: \;
  r_1 \xrightarrow{\text{hb}} r_2 \implies
  \text{vis}(r_1) \subseteq \text{vis}(r_2)
\end{equation}
where $\xrightarrow{\text{hb}}$ denotes the happens-before relation
and $\text{vis}(r)$ is the set of operations visible to request $r$.

\subsection{Fault Tolerance}

Describe how \system handles failures:

\para{Node Failures}
[How the system handles crashed or slow nodes.]

\para{Network Partitions}
[How the system handles network partitions.]

\para{Recovery}
[How the system recovers after failures.]

%----------------------------------------------------------------------
\section{Correctness}
\label{sec:correctness}

% SOSP papers in areas like distributed systems, storage, and
% concurrency often include formal correctness arguments.

We prove that \system maintains [property] under [failure model].

\begin{theorem}
  \label{thm:safety}
  Under the failure model of \cref{sec:background}, \system
  guarantees [safety property]: for all executions $E$,
  [formal statement].
\end{theorem}

\begin{proof}[Proof sketch]
  By induction on the number of operations. The base case holds
  because [reason]. For the inductive step, [key argument].
  Full proof in the supplementary material.
\end{proof}

%----------------------------------------------------------------------
\section{Implementation}
\label{sec:implementation}

We implement \system in approximately [X]K lines of [language].
Key implementation details include:

\para{Storage Layer}
[Describe the storage implementation.]

\para{Network Layer}
[Describe the networking implementation.]

\para{Concurrency Control}
[Describe the concurrency control mechanism.]

% Example code snippet (common in SOSP papers)
\Cref{lst:api} shows the client API for \system.

\begin{figure}[t]
\begin{lstlisting}[caption={Client API for \system. The interface
  provides [property] guarantees.}, label={lst:api},
  language=Python]
class Client:
  def get(self, key, consistency="strong"):
    """Read with configurable consistency."""
    ts = self._get_timestamp()
    return self._send_read(key, ts,
                           consistency)

  def put(self, key, value):
    """Write with durability guarantee."""
    ts = self._get_timestamp()
    ack = self._send_write(key, value, ts)
    return ack.committed
\end{lstlisting}
\end{figure}

%----------------------------------------------------------------------
\section{Evaluation}
\label{sec:evaluation}

We evaluate \system to answer:
\begin{enumerate}
  \item How does \system compare to state-of-the-art systems?
  \item What is the cost of [guarantee] in terms of performance?
  \item How does \system perform under failures?
  \item What is the contribution of each mechanism?
\end{enumerate}

\subsection{Experimental Setup}
\label{sec:eval:setup}

\para{Testbed}
We run experiments on [N] machines in [cloud/cluster], each with
[CPU], [X]\GB RAM, [Y]\GB SSD, connected via [network].

\para{Workloads}
We use [standard benchmarks] and [production traces].
\Cref{tab:workloads} summarizes the workload characteristics.

\para{Baselines}
We compare against:
(1)~[System A]~\cite{ghemawat2003gfs},
(2)~[System B]~\cite{corbett2013spanner}, and
(3)~[System C]~\cite{decandia2007dynamo}.

\begin{table}[t]
  \caption{Workload characteristics. Workloads span different
    read/write ratios and access patterns.}
  \label{tab:workloads}
  \centering
  \begin{small}
  \begin{tabular}{@{}lrrcl@{}}
    \toprule
    \textbf{Workload} & \textbf{Ops/s} & \textbf{Data} &
    \textbf{R:W} & \textbf{Pattern} \\
    \midrule
    YCSB-A   & 100K  & 10\,GB  & 50:50  & Uniform   \\
    YCSB-B   & 100K  & 10\,GB  & 95:5   & Zipfian   \\
    YCSB-C   & 100K  & 10\,GB  & 100:0  & Zipfian   \\
    YCSB-F   & 50K   & 10\,GB  & 50:50  & RMW       \\
    Production& 200K  & 100\,GB & 80:20  & Zipfian   \\
    \bottomrule
  \end{tabular}
  \end{small}
\end{table}

\subsection{End-to-End Performance}
\label{sec:eval:e2e}

\Cref{tab:e2e} shows the end-to-end performance comparison.
\system achieves [X]\% higher throughput and [Y]\% lower tail
latency compared to [best baseline].

\begin{table}[t]
  \caption{End-to-end performance comparison across workloads.
    Bold indicates best result.}
  \label{tab:e2e}
  \centering
  \begin{small}
  \begin{tabular}{@{}lcccc@{}}
    \toprule
    & \multicolumn{2}{c}{\textbf{YCSB-A}} &
      \multicolumn{2}{c}{\textbf{Production}} \\
    \cmidrule(lr){2-3} \cmidrule(lr){4-5}
    \textbf{System} & \textbf{Kops/s} & \textbf{p99 (ms)} &
    \textbf{Kops/s} & \textbf{p99 (ms)} \\
    \midrule
    System A         & 85.2   & 12.4  & 142.1  & 18.7 \\
    System B         & 72.1   & 15.8  & 128.4  & 22.1 \\
    System C         & 98.4   & 8.2   & 165.3  & 11.5 \\
    \textbf{\system} & \textbf{124.6} & \textbf{5.1} &
                       \textbf{201.8} & \textbf{7.3} \\
    \bottomrule
  \end{tabular}
  \end{small}
\end{table}

\subsection{Performance Under Failures}
\label{sec:eval:failure}

\begin{figure}[t]
  \centering
  \fbox{\parbox{0.9\columnwidth}{\centering\vspace{3em}
    \textit{Time-series: throughput and latency during node failure \\
    and recovery, showing impact duration and recovery time}
  \vspace{3em}}}
  \caption{Performance during node failure at $t=60$s. \system
    recovers within [X]\ms with [Y]\% throughput drop, compared
    to [Z]\ms and [W]\% drop for [baseline].}
  \label{fig:failure}
\end{figure}

\subsection{Microbenchmarks}
\label{sec:eval:micro}

We isolate the performance of key mechanisms.

\begin{figure}[t]
  \centering
  \fbox{\parbox{0.9\columnwidth}{\centering\vspace{3em}
    \textit{Latency breakdown or CDF comparing mechanisms}
  \vspace{3em}}}
  \caption{Latency CDF for individual operations. \system's
    [mechanism] adds only [X]\us overhead to the critical path.}
  \label{fig:microbench}
\end{figure}

\subsection{Ablation Study}
\label{sec:eval:ablation}

\Cref{fig:ablation} shows the contribution of each mechanism.

\begin{figure}[t]
  \centering
  \fbox{\parbox{0.9\columnwidth}{\centering\vspace{3em}
    \textit{Bar chart: \system variants with mechanisms disabled}
  \vspace{3em}}}
  \caption{Ablation study on YCSB-A. Mechanism A contributes
    [X]\% and Mechanism B contributes [Y]\% of the improvement.}
  \label{fig:ablation}
\end{figure}

%----------------------------------------------------------------------
\section{Discussion}
\label{sec:discussion}

\para{Lessons Learned}
[Share insights from designing and building the system.]

\para{Limitations}
[Honest discussion of where \system falls short.]

\para{Applicability}
[Discuss how the principles generalize to other systems.]

%----------------------------------------------------------------------
\section{Related Work}
\label{sec:related}

% Organize by theme. SOSP reviewers expect thorough related work.

\para{Distributed Storage Systems}
GFS~\cite{ghemawat2003gfs}, Dynamo~\cite{decandia2007dynamo}, and
Spanner~\cite{corbett2013spanner} address [aspect]. \system differs by
[distinction].

\para{Consensus and Coordination}
Paxos~\cite{lamport1998paxos} and ZooKeeper~\cite{hunt2010zookeeper}
provide [guarantee]. \system builds on these foundations while
[extending/relaxing] [aspect].

\para{[Other Category]}
[Other related work] explores [approach]. \system complements these
by [distinction].

%----------------------------------------------------------------------
\section{Conclusion}
\label{sec:conclusion}

We presented \system, a [type of system] that [key capability].
Through [mechanisms], \system achieves [X]$\times$ improvement in
[metric] while guaranteeing [property].
Our experience building \system reveals [key lesson], suggesting
[future direction] for systems research.

%----------------------------------------------------------------------
% Acknowledgments (only in camera-ready, remove for submission)
% \begin{acks}
% We thank the anonymous reviewers and our shepherd for their
% invaluable feedback. This work was supported by [funding].
% \end{acks}

\bibliographystyle{ACM-Reference-Format}
\bibliography{references}

%----------------------------------------------------------------------
% ARTIFACT EVALUATION (optional, after acceptance)
% SOSP offers optional artifact evaluation. If your paper is accepted,
% consider preparing your artifact for evaluation.
% See: https://sysartifacts.github.io/

\end{document}
